\documentclass[12pt,titlepage]{article}

\usepackage{amsmath}
\usepackage{xcolor}
\usepackage{mathtools}
\usepackage{indentfirst}
\usepackage{hhline}

\title{\bf Largest $m$-cube in an $n$-cube: \\
       Examples of Solving Specific Cases
      }

\author{Greg Huber \\
        Terry J. Ligocki
       }

\begin{document}

\maketitle

\begin{abstract}
Here are examples of solving for the largest m-cube that can be fit in an
n-cube but using empirical data and the problem constraints to produce
algegraic equations that have a root that is the solution to the given
problem.  The examples are for (m,n) equal to (1,2), (2,3), and (3,4).
\end{abstract}

\section{Introduction}

We will be using two different matrix representations of the m-cube.  One
version, V, will have the vertices of the m-cube in n-dimensional space as
rows of the matrix.  Thus, this matrix will be $2^m$ x $n$.

The other matrix, E, will contain the edge vectors of the m-cube as rows.
These can be computed from all pairs of adjacent vertices but many of these
vectors will be identical since there are many parallel edges.  Thus, this
matrix will be $m$ x $n$.

The form of the vertex matrix will be estimated from empirical results.
Specifically, the best empirical result will be used to determine which
coordinates are probably $\pm 1/2$ and which coordinates are equal to other
coordinates.

\section{(1,2) Case}

Here the vertex matrix is:
\begin{equation*}
V =
\begin{bmatrix*}[r]
v_0 \\
v_1
\end{bmatrix*}
=
\begin{bmatrix*}[r]
-1/2 & -1/2  \\
 1/2 &  1/2 
\end{bmatrix*}
\end{equation*}
And the edge matrix is:
\begin{equation*}
E =
\begin{bmatrix*}[r]
e_0
\end{bmatrix*}
=
\begin{bmatrix*}[r]
v_1 - v_0
\end{bmatrix*}
=
\begin{bmatrix*}[r]
 1 & 1
\end{bmatrix*}
\end{equation*}
Since there is only one edge there are no constraints beyond this and we know
that f(1,2) = $|e_0|$ = $\sqrt{2}$.

\section{(2,3) Case}

Here the vertex matrix is:
\begin{equation*}
V =
\begin{bmatrix*}[r]
v_{0,0} \\
v_{1,0} \\
v_{0,1} \\
v_{1,1}
\end{bmatrix*}
=
\begin{bmatrix*}[r]
-1/2 & -1/2 &   -x \\
 1/2 &   -x & -1/2 \\
-1/2 &    x &  1/2 \\
 1/2 &  1/2 &    x \\
\end{bmatrix*}
\end{equation*}
And the edge matrix is:
\begin{equation*}
E =
\begin{bmatrix*}[r]
e_{1,0} \\
e_{0,1}
\end{bmatrix*}
=
\begin{bmatrix*}[r]
v_{1,0} - v_{0,0} \\
v_{0,1} - v_{0,0}
\end{bmatrix*}
=
\begin{bmatrix*}[r]
v_{1,1} - v_{0,1} \\
v_{1,1} - v_{1,0}
\end{bmatrix*}
=
\begin{bmatrix*}[r]
1 & -x + 1/2 & x - 1/2 \\
0 &  x + 1/2 & x + 1/2
\end{bmatrix*}
\end{equation*}
The constraints are:
\begin{align*}
e_{1,0} \cdot e_{1,0} &= |e_{1,0}|^2 = f(2,3)^2 \\
e_{1,0} \cdot e_{0,1} &= 0                      \\
e_{0,1} \cdot e_{0,1} &= |e_{0,1}|^2 = f(2,3)^2
\end{align*}
Substituting for $e_{1,0}$ and $e_{0,1}$:
\begin{align*}
1 + (-x + 1/2)^2 + (x - 1/2)^2           &= f(2,3)^2 \\
(-x + 1/2)(x + 1/2) + (x - 1/2)(x + 1/2) &= 0        \\
0 + (x + 1/2)^2 + (x + 1/2)^2            &= f(2,3)^2
\end{align*}
The second equation is satisfied for all $x$.  The first and third equations
combine to give:
\begin{align*}
1 + 2 (x - 1/2)^2             &= 2 (x + 1/2)^2 \\
2 (x + 1/2)^2 - 2 (x - 1/2)^2 &= 1             \\
4 x                           &= 1             \\
x                             &= 1/4
\end{align*}
So:
\begin{equation*}
f(2,3) = \sqrt{2 (x + 1/2)^2} = \sqrt{9/8}
\end{equation*}

\section{(3,4) Case}

Here the vertex matrix is:
\begin{equation*}
V =
\begin{bmatrix*}[r]
v_{0,0,0} \\
v_{1,0,0} \\
v_{0,1,0} \\
v_{1,1,0} \\
v_{0,0,1} \\
v_{1,0,1} \\
v_{0,1,1} \\
v_{1,1,1}
\end{bmatrix*}
=
\begin{bmatrix*}[r]
   x &  1/2 &    y &  1/2 \\
 1/2 &  1/2 &    z &   -w \\
   z & -1/2 &  1/2 &    w \\
   y & -1/2 &    x & -1/2 \\
  -y &  1/2 &   -x &  1/2 \\
  -z &  1/2 & -1/2 &   -w \\
-1/2 & -1/2 &   -z &    w \\
  -x & -1/2 &   -y & -1/2
\end{bmatrix*}
\end{equation*}
And the edge matrix is:
\begin{align*}
E =
\begin{bmatrix*}[r]
e_{1,0,0} \\
e_{0,1,0} \\
e_{0,0,1}
\end{bmatrix*}
&=
\begin{bmatrix*}[r]
v_{1,0,0} - v_{0,0,0} \\
v_{0,1,0} - v_{0,0,0} \\
v_{0,0,1} - v_{0,0,0}
\end{bmatrix*}
= 
\begin{bmatrix*}[r]
v_{1,1,0} - v_{0,1,0} \\
v_{0,1,1} - v_{0,0,1} \\
v_{1,0,1} - v_{1,0,0}
\end{bmatrix*}
= 
\begin{bmatrix*}[r]
v_{1,0,1} - v_{0,0,1} \\
v_{1,1,0} - v_{1,0,0} \\
v_{0,1,1} - v_{0,1,0}
\end{bmatrix*}
= 
\begin{bmatrix*}[r]
v_{1,1,1} - v_{0,1,1} \\
v_{1,1,1} - v_{1,0,1} \\
v_{1,1,1} - v_{1,1,0}
\end{bmatrix*}
\\
&
\\
&=  
\begin{bmatrix*}[r]
-x + 1/2 &  0 & -y +   z & -w - 1/2 \\
-x +   z & -1 & -y + 1/2 &  w - 1/2 \\
-x -   y &  0 & -y -   x &        0
\end{bmatrix*}
\\
&
\\
&=  
\begin{bmatrix*}[r]
a        &  0 & b        &  c   \\
d        & -1 & a+b-d    & -c-1 \\
2a+b-d-1 &  0 & 2a+b-d-1 &  0   \\
\end{bmatrix*}
\end{align*}
where $a = -x+1/2$, $b = -y+z$, $c = -w+1/2$, and $d = -x+z$.
The constraints are:
\begin{align*}
e_{1,0,0} \cdot e_{1,0,0} &= |e_{1,0,0}|^2 = f(3,4)^2 \\
e_{1,0,0} \cdot e_{0,1,0} &= 0                        \\
e_{1,0,0} \cdot e_{0,0,1} &= 0                        \\
e_{0,1,0} \cdot e_{0,1,0} &= |e_{0,1,0}|^2 = f(3,4)^2 \\
e_{0,1,0} \cdot e_{0,0,1} &= 0                        \\
e_{0,0,1} \cdot e_{0,0,1} &= |e_{0,0,1}|^2 = f(3,4)^2
\end{align*}
Substituting for $e_{1,0,0}$, $e_{0,1,0}$, and $e_{0,0,1}$:
\begin{align}
a^2 + b^2 + c^2               &= f(3,4)^2 \label{con1:1} \\
ad + b(a+b-d) - c(c+1)        &= 0        \label{con1:2} \\
(a+b)(2a+b-d-1)               &= 0        \label{con1:3} \\
d^2 + 1 + (a+b-d)^2 + (c+1)^2 &= f(3,4)^2 \label{con1:4} \\
(a+b)(2a+b-d-1)               &= 0        \label{con1:5} \\
2(2a+b-d-1)^2                 &= f(3,4)^2 \label{con1:6}
\end{align}
Constraint (\ref{con1:3}) and (\ref{con1:5}) are the same.  Further, since
$(2a+b-1) \ne 0$, $b = -a$.  So the constraints become:
\begin{align*}
2a^2 + c^2         &= f(3,4)^2  \\
2ad - c(c+1)       &= 0         \\
2d^2 + (c+1)^2 + 1 &= f(3,4)^2  \\
2(a-d-1)^2         &= f(3,4)^2 
\end{align*}
Letting $m = f(3,4)^2$, these equations can be reduced to:
\begin{equation*}
16 m^8 - 224 m^7 + 984 m^6 - 1528 m^5 +625 m^4 + 352 m^3 + 32 m^2 - 512 m +
256 = 0
\end{equation*}
and we are interested in the smallest, positive, real root.  This is:
\begin{align*}
m &= \frac{1}{12}\left
     (21 + \sqrt{3c} - \sqrt{1449 + 6480\sqrt{\frac{3}{c}} - 3c}\right)
   = 1.0149247893...
\\
c &= 161 + (184841 - 4320\sqrt{1290})^{1/3}
         + (184841 + 4320\sqrt{1290})^{1/3}
\end{align*}
and this means $l = 1.0074347568...$

\end{document}
